\section{Objetivos}
\subsection{Objetivo general}
\begin{itemize}

    \item Diseñar e implementar una estrategia eficiente de gestión de memoria en simuladores cuánticos basada en técnicas de compresión sin pérdida de precisión, con el fin de optimizar el uso de recursos computacionales sin comprometer la exactitud numérica de las simulaciones.

\end{itemize}

\subsection{Objetivos específicos}

\begin{itemize}
    \item Analizar el impacto del uso de herramientas de compresión de datos sin pérdida de precisión en el consumo de memoria, el tiempo de ejecución y la verificación de igualdad exacta de los resultados respecto a una referencia no comprimida.
    \item Diseñar una estrategia integral de gestión de memoria basada en compresión sin pérdida de precisión, definiendo el algoritmo, los parámetros de compresión‑descompresión y los puntos de integración con el simulador cuántico seleccionado.
    \item Desarrollar e integrar la estrategia diseñada en el simulador elegido, optimizando la sobrecarga computacional para garantizar un ahorro significativo de memoria sin afectar la exactitud numérica de la simulación.
    \item Implementar dos algoritmos cuánticos distintos en el simulador cuántico elegido para evaluar el uso de memoria, velocidad y exactitud con y sin compresión de datos.
    \item Comparar los resultados obtenidos con estrategias tradicionales y métodos que emplean compresión con pérdida de precisión, identificando ventajas y desventajas de cada enfoque.
\end{itemize}
